\section{Introduction}

The successful implementation of the \textbf{AgriGov Market} platform depends on a robust architectural foundation that translates the functional requirements identified in the previous chapter into a concrete technical reality. Transitioning from the "what" to the "how" requires a rigorous design process to ensure that the system is scalable, secure, and capable of handling the complex interactions between various agricultural actors. By defining the internal structure of the software and its data organization, we move away from abstract needs toward a structured blueprint that guides the actual development and deployment of the platform \cite{Pressman2024}.

The primary objective of this chapter is to detail the design of the system's architecture and its technical components. We begin by presenting the final design class diagram, which establishes the logical structure of the application and the relationships between its core entities. To visualize the dynamic flow of the platform, a global activity diagram is provided, illustrating how processes such as order placement and logistics tracking are executed. Furthermore, we describe the transformation of the object-oriented model into a relational database schema, ensuring data integrity and efficient retrieval for the Ministry's monitoring tools \cite{Bass2021}.

The chapter is structured as follows. First, we present the design class diagram and the global activity flow. We then detail the database design process, including the mapping rules and the resulting relational schema. Subsequently, we showcase the final user interfaces to demonstrate the system's usability. Finally, we discuss the technical architecture, the deployment strategies, and the specific technology stack—including frameworks and development environments—used to bring the AgriGov Market platform to fruition.